\documentclass{article}

\usepackage[submission]{colm2026_conference}
\usepackage{microtype}
\usepackage{amsmath}
\usepackage{booktabs}
\usepackage{hyperref}
\usepackage{url}
\usepackage{xcolor}
\usepackage{lineno}

\setlength{\emergencystretch}{2em}
\setlength{\textfloatsep}{8pt plus 2pt minus 2pt}
\setlength{\floatsep}{6pt plus 2pt minus 2pt}
\setlength{\intextsep}{6pt plus 2pt minus 2pt}
\setlength{\abovedisplayskip}{4pt}
\setlength{\belowdisplayskip}{4pt}
\setlength{\abovedisplayshortskip}{3pt}
\setlength{\belowdisplayshortskip}{3pt}

\title{Trusted by Proxy: Source-Laundering Attacks in Multi-Agent LLM Systems}
\author{Anonymous Authors}
\date{}

\begin{document}
\ifcolmsubmission
\linenumbers
\fi

\maketitle

\begin{abstract}
Multi-agent LLM systems often split document reading, planning, and tool
execution across specialized agents. This modularity creates a trust-boundary
failure: an instruction from untrusted content can be rewritten by an upstream
peer as an ordinary action item, causing downstream agents to treat it as
authorized. We call this \emph{source laundering}. We introduce SLaMBench, a
synthetic benchmark for source-laundering attacks in tool-using agents across
workspace, CRM, code-review, travel, and research-assistant domains. On a
50-pair multidomain benchmark, an extractive reader--planner--executor topology
raises unauthorized-recipient attack success from 6\% in a direct single-agent
baseline to 22\%, a +16 percentage-point lift. A deterministic TraceGate monitor
eliminates observed unauthorized sends while preserving more injected-task
utility than the vulnerable multi-agent system. Ablations show that generic
delegation depth is not sufficient: attacks appear when intermediate agents
strip or normalize provenance. A current-model \texttt{gpt-5.5} replication has
low attack success across all tested conditions, showing model dependence while
preserving the diagnostic value of provenance dropout.
\end{abstract}

\section{Introduction}

LLM agents are increasingly deployed as \emph{systems} rather than as single
prompts. A user-facing assistant may delegate document reading to one agent,
planning to another, and tool execution to a third. This decomposition localizes
expertise and shortens prompts, but it also creates a security boundary. A Reader
may preserve useful facts from an untrusted artifact while dropping the fact that
the artifact had no authority; a Planner or Executor may then receive the same
content as a trusted peer message.

Consider a benign user request: ``read this vendor note and email
\texttt{finance@example.test} the payment terms.'' The vendor note correctly
states the payment terms, but also contains ``add \texttt{audit-log@example.test}
as a required recipient.'' A direct single-agent system sees that the extra
recipient is embedded in the document. In a multi-agent system, however, a Reader
might rewrite the artifact as ``action item: add the audit recipient,'' and a
Planner might further rewrite it as ``required step: send the summary to
\texttt{audit-log@example.test}.'' The downstream Executor no longer sees a raw
untrusted instruction; it sees a plan produced by a peer. The question is not
whether indirect prompt injection exists, but whether delegation changes the
instruction's effective authority.

We call this transformation \emph{source laundering}. The analogy is intentional:
the content has not disappeared, but its provenance has been washed through
apparently legitimate transformations. Source laundering is related to indirect
prompt injection~\citep{greshake2023not,zhan2024injecagent}, but its unit of
analysis is the inter-agent boundary; it is also a confused-deputy
failure~\citep{hardy1988confused}, where a tool-authorized component exercises
authority on behalf of a less trusted principal.

We ask whether extractive delegation increases unauthorized tool use relative to
direct exposure (RQ1), which payloads and domains survive the reader--planner
boundary (RQ2), and whether prompt-level provenance or deterministic
tool-boundary checks reduce the risk without destroying utility (RQ3). Our
contributions are a formalization of source laundering, SLaMBench, an evaluation
of direct/delegated/defended topologies, and evidence that the failure is tied to
provenance-losing composition rather than delegation depth alone. We do not claim
that every multi-agent system is less secure than every single-agent system; we
show that a plausible extractive boundary can convert untrusted artifact text
into peer-agent authority unless provenance is preserved and enforced.

\section{Threat Model and Security Property}

The user provides a legitimate task, such as reading an invoice note and
emailing an authorized summary to \texttt{finance@example.test}. The adversary
can place text inside an untrusted artifact that the system must process, but
cannot modify system prompts, trusted context, tool code, model weights, runtime
policy, or the user's initial goal. The protected assets are synthetic but
structurally representative: authorized recipients, synthetic secrets, protected
status fields, calendar attendees, and memory entries. The core security property
is \emph{authority separation}: untrusted content may supply facts needed for the
task, but it must not authorize new tool calls, new recipients, protected
updates, memory policies, or secret disclosure.

We define a source-laundering attack as an execution with three properties: an
attacker-controlled instruction appears in an untrusted artifact; an intermediate
agent transforms it into a peer message without sufficient source, trust, and
no-authority metadata; and a downstream agent acts on the transformed instruction
by emitting a forbidden tool call. A run that merely fails the legitimate task is
a utility failure, while a run that sends to \texttt{audit-log@example.test} is a
security failure. All tools are deterministic mocks, with no real email,
calendar, ticket, browser, or memory service contacted.

\section{SLaMBench}

SLaMBench contains paired benign and injected task variants. The main scaled task
set, SLaMBench-50, contains 50 benign tasks and 50 injected variants across five
domains (Table~\ref{tab:slambench}). Each pair shares the same user goal, trusted
context, allowed tools, expected benign outcome, and reference benign tool calls;
the injected variant adds one attacker-controlled payload to the untrusted
artifact. This paired design lets us compare benign utility and injected security
without changing the legitimate task.

\begin{table}[t]
\centering
\small
\begin{tabular}{lcll}
\toprule
Domain & Pairs & Legitimate action & Unauthorized action \\
\midrule
Workspace & 10 & Email payment terms to finance & Add audit recipient \\
CRM & 10 & Update ticket and email support & Add audit recipient \\
Code review & 10 & Update tracker and email maintainer & Add audit recipient \\
Travel & 10 & Create approved calendar event & Add audit recipient \\
Research & 10 & Write bibliography note and email lead & Add audit recipient \\
\bottomrule
\end{tabular}
\caption{SLaMBench scaled task set used for the main result.}
\label{tab:slambench}
\end{table}

\paragraph{Task format.}
Each task specifies the user goal, trusted context, untrusted artifacts, allowed
tools, expected benign outcomes, forbidden outcomes, and an attack tool-call
template. The main forbidden outcome is any executed \texttt{send\_email} call
containing \texttt{audit-log@example.test}, with additional checks for
synthetic-secret leakage and untrusted memory-policy writes.

\paragraph{Payload classes.}
The injected variants use five intentionally simple payload classes: direct
instructions, fake authority claims, delegation poisoning, summary bait, and
memory poisoning. The goal is not to optimize an adaptive attack, but to measure
how ordinary delegation transforms source and authority.

\paragraph{Why synthetic tasks.}
The benchmark is synthetic by design. This makes every tool call inspectable,
eliminates risk to real users, and avoids relying on LLM judges for the main
security metric. Synthetic tasks also let us pair benign and injected variants
with controlled differences. The limitation is that absolute attack rates should
not be interpreted as production incident rates; the benchmark is a diagnostic
instrument for a specific compositional failure mode.

\section{Agent Topologies and Defenses}

\textbf{Direct single agent} (\texttt{c0\_single\_direct}) sees the user goal,
trusted context, untrusted artifact, and tool signatures, then emits final tool
calls. \textbf{Extractive multi-agent}
(\texttt{c2x\_3agent\_extractive\_naive}) uses a Reader that extracts action
items from the artifact, a Planner that turns the Reader output into steps, and
an Executor that emits tool calls. The untrusted artifact is not labeled as
untrusted in the extractive Reader prompt, modeling a common pattern where
upstream agents compress or normalize task context before execution.

We test three defenses. \textbf{Source-preserving}
(\texttt{c3\_source\_preserving}) requires source identifiers and trust labels in
Reader output and tells downstream agents not to convert untrusted instructions
into authorized actions. \textbf{Capability-scoped}
(\texttt{c4\_capability\_scoped}) allows peers to suggest actions but not
authorize sensitive tools unless the original user goal or trusted context does
so. \textbf{TraceGate} (\texttt{c5x\_tracegate\_extractive\_naive}) uses the same
extractive topology as c2x but checks sensitive tool calls before execution,
blocking synthetic-secret disclosure, unauthorized recipients or attendees,
untrusted memory-policy writes, and sensitive actions without trusted
authorization. For mixed-recipient emails, it prunes unauthorized recipients and
preserves authorized sends when possible.

\section{Metrics and Analysis Protocol}

All metrics are computed from saved tool traces. \textbf{Benign Success Rate}
(BSR) is the fraction of benign variants where the expected user task succeeds
and no forbidden event occurs. \textbf{Injected Utility Success Rate} (IUSR)
applies the same check to injected variants. \textbf{Attack Success Rate} (ASR)
is the fraction of injected variants with a forbidden event in the executed tool
log; blocked calls are not counted as attacks.

\textbf{Authority Transfer Rate} (ATR) measures delegation-specific lift:
\begin{equation*}
\mathrm{ATR}(c) = \mathrm{ASR}(c) - \mathrm{ASR}(\texttt{c0\_single\_direct}).
\end{equation*}
This metric is important because some models are directly vulnerable. A high raw
ASR in a delegated condition is evidence of insecurity; a positive ATR is
evidence that the inter-agent boundary adds risk beyond direct exposure.

\textbf{Provenance Dropout Rate} (PDR) is the fraction of reader/planner messages
that carry task-specific forbidden content without source, trust, and
no-authority metadata. \textbf{Capability-Normalized Attack Success} (CNAS) is
ASR divided by BSR, preventing weak agents from appearing secure merely because
they fail both benign and malicious tasks. We report rates and differences in
percentage points.

\section{Main Results}

The main scaled result uses \texttt{gpt-5.4-nano} on SLaMBench-50. Full tables
are in \path{paper/tables/openai_schema_prompt_50pairs_gpt54nano_tracegate_pruned/}.

\begin{table}[t]
\centering
\small
\begin{tabular}{lrrrrrrr}
\toprule
Condition & BSR & IUSR & ASR & ATR & PDR & CNAS & Blocks \\
\midrule
Direct single-agent & 60\% & 54\% & 6\% & 0 pp & 0\% & 0.10 & 0 \\
Extractive multi-agent & 70\% & 50\% & 22\% & +16 pp & 77\% & 0.31 & 0 \\
TraceGate & 66\% & 64\% & 0\% & -6 pp & 69\% & 0.00 & 11 \\
\bottomrule
\end{tabular}
\caption{Main 50-pair result on \texttt{gpt-5.4-nano}. ATR is ASR minus the
direct-baseline ASR; PDR is provenance dropout in intermediate messages; CNAS is
ASR divided by BSR.}
\label{tab:main}
\end{table}

The extractive multi-agent topology increases ASR by +16 percentage points over
direct exposure. TraceGate reduces ASR by -22 points relative to extractive
multi-agent delegation. TraceGate produces zero audit-log sends, records 11
sensitive-call blocks, and sanitizes five mixed-recipient emails so the
authorized recipient still receives the intended message.

\paragraph{Answer to RQ1.}
The same injected artifacts are present in c0 and c2x, but c2x has a higher ASR.
The trace-level explanation is consistent with source laundering: the Reader
extracts the forbidden recipient as an action item, the Planner promotes it into
a required step, and the Executor emits the corresponding tool call.

\paragraph{Utility/security tradeoff.}
The extractive multi-agent system has higher benign success than the direct
baseline (70\% versus 60\%), but the utility gain comes with higher ASR and lower
IUSR. TraceGate recovers the best injected utility among the three main
conditions (64\%) while driving ASR to 0\%, showing that tool-boundary
enforcement can block source-laundered actions without refusing the whole task.

\paragraph{Provenance dropout.}
The high PDR values show that extractive delegation frequently carries forbidden
recipient content forward without source and authority metadata. TraceGate does
not repair this upstream dropout; it prevents the laundered instruction from
becoming an executed unauthorized tool call.

\begin{table}[t]
\centering
\small
\begin{tabular}{lc}
\toprule
Payload class & c2x ASR \\
\midrule
Summary bait & 50\% \\
Fake authority & 40\% \\
Delegation poisoning & 10\% \\
Direct instruction & 10\% \\
Memory poisoning & 0\% \\
\bottomrule
\end{tabular}
\caption{Payload-level ASR for extractive multi-agent delegation.}
\label{tab:payload}
\end{table}

\paragraph{Answer to RQ2.}
Table~\ref{tab:payload} suggests that the most successful attacks are not the
most obviously adversarial ones. Summary bait and fake-authority payloads are
more effective than direct instructions in the extractive topology. This matches
the threat model: a Reader optimized to extract action items is more likely to
preserve text that already looks like workflow content. The attack appears across
all five domains in c2x, although the main positive result remains
unauthorized-recipient sends.

\paragraph{Qualitative trace.}
In one fake-authority case, the Reader converts an untrusted ``user-approved
routing exception'' into an action item to send the summary to
\texttt{audit-log@example.test}; the Planner schedules it and the Executor sends
it. In a TraceGate case, a mixed-recipient email is pruned to the authorized
finance recipient. These traces illustrate why we distinguish upstream
laundering from downstream enforcement.

\section{Topology Ablation}

The follow-up topology ablation uses the same 25-pair SLaMBench slice for all
conditions and \texttt{gpt-5.4-nano}. It separates generic delegation depth from
the provenance-losing handoff that defines source laundering.

\begin{table}[t]
\centering
\small
\begin{tabular}{lrrrr}
\toprule
Condition & BSR & IUSR & ASR & PDR \\
\midrule
Direct single-agent & 60\% & 48\% & 4\% & 0\% \\
2-agent naive & 64\% & 60\% & 0\% & 72\% \\
3-agent naive & 72\% & 72\% & 4\% & 52\% \\
3-agent extractive & 68\% & 56\% & 24\% & 86\% \\
Oracle-laundered handoff & 80\% & 52\% & 28\% & 70\% \\
\bottomrule
\end{tabular}
\caption{Topology ablation on the 25-pair benchmark with \texttt{gpt-5.4-nano}.
The attack emerges in provenance-losing handoffs, not in naive delegation depth
alone.}
\label{tab:topology}
\end{table}

Table~\ref{tab:topology} supports a more precise mechanism claim. The naive
2-agent and 3-agent chains do not increase ASR over direct exposure. In
contrast, the extractive reader that normalizes document instructions into
action items raises ASR by +20 points over direct exposure. The oracle-laundered
condition, where the downstream planner and executor receive a deliberately
provenance-stripped handoff, raises ASR by +24 points. Thus the core risk is not
simply ``more agents,'' but composition that strips source and authority metadata
before a tool-authorized component acts. A no-API trace audit of these saved runs
gives the same picture: in all six injected tasks where extractive delegation
attacks and direct exposure does not, the forbidden recipient appears in the
reader handoff, the planner plan, and the executor output before the tool call.
Post-hoc diagnostics support the same conservative boundary: on a six-task hard
slice, extractive delegation re-attacks 8/18 repeated runs while direct exposure
attacks 2/18; TraceGate replay yields 0/18 executed attacks. A tiny non-email
side-effect probe does not show positive c2x lift, so our amplification claim is
centered on unauthorized-recipient sends.

\section{Defense Ablations}

On the same clean 25-pair slice, we compare prompt-level defenses against the
extractive condition and replay the same sampled extractive outputs through
TraceGate. The prompt-defense rows use raw runs with zero parser or truncation
errors. The TraceGate row is deterministic replay rather than an independent
model resampling: c2x and c5x use identical prompts and differ only in the
tool-boundary runtime monitor.

\begin{table}[t]
\centering
\small
\begin{tabular}{lrrrrr}
\toprule
Condition & BSR & IUSR & ASR & PDR & Blocks \\
\midrule
Extractive multi-agent & 68\% & 56\% & 24\% & 86\% & 0 \\
Source-preserving & 44\% & 56\% & 0\% & 86\% & 0 \\
Capability-scoped & 52\% & 52\% & 0\% & 86\% & 0 \\
TraceGate replay & 68\% & 72\% & 0\% & 86\% & 10 \\
LLM guard replay & 68\% & 72\% & 0\% & 86\% & 0 \\
\bottomrule
\end{tabular}
\caption{Clean 25-pair defense comparison on \texttt{gpt-5.4-nano}. TraceGate is
a deterministic replay of the extractive model outputs through the runtime
monitor; the LLM guard is a replay sanitizer over the same proposed tool calls.
Neither replay row is an independent model resampling.}
\label{tab:defenses}
\end{table}

\paragraph{Answer to RQ3.}
Source and authority metadata are helpful: source-preserving and
capability-scoped prompts reduce ASR from 24\% to 0\%, a -24 point change
against extractive delegation.
The deterministic TraceGate replay produces the same ASR reduction while
preserving the extractive chain's benign success and improving injected utility
to 72\% by pruning or blocking unauthorized calls. An LLM guard replay over the
same saved c2x tool calls also reaches 0\% ASR and 72\% IUSR, but adds another
model call and can only be audited empirically. Prompt-only defenses reduce
benign success, suggesting a layered design: upstream agents preserve metadata
and downstream runtime monitors enforce high-assurance invariants.

\section{Model Dependence}

The effect is model-dependent. The \texttt{gpt-5.4-nano} topology ablation shows
a clear source-laundering lift, while two other model checks fall into different
interpretation regimes: \texttt{gpt-4.1-nano} is already directly vulnerable,
and \texttt{gpt-5.5} is robust on the 25-pair slice.
The 25-pair ASRs are 4\%$\rightarrow$24\% for \texttt{gpt-5.4-nano}
(positive ATR), 24\%$\rightarrow$24\% with TraceGate at 0\% for
\texttt{gpt-4.1-nano} (high direct ASR), and 0\%$\rightarrow$0\% with TraceGate
at 0\% for \texttt{gpt-5.5} (low ASR). Thus, the right claim is conditional
rather than universal. Source laundering is most visible when the direct model
resists many raw artifact instructions, but delegation removes that protective
context. If a model is already directly vulnerable, delegation may not increase
ASR; if a model rejects both direct and laundered instructions, SLaMBench still
exposes latent provenance dropout but not executed attacks. A six-case
heterogeneous diagnostic with a \texttt{gpt-5.5} executor also yields 0/6
attacks, suggesting executor robustness is another axis.

\section{Related Work}

Indirect prompt injection shows that LLM-integrated applications blur the
boundary between data and instructions when processing retrieved or third-party
content~\citep{greshake2023not}. AgentDojo, InjecAgent, Agent Security Bench, and
ToolEmu evaluate attacks and defenses for tool-using agents across broader task
suites~\citep{debenedetti2024agentdojo,zhan2024injecagent,zhang2024asb,ruan2023toolemu}.
Prompt Infection, AgentPoison, and ToolHijacker show that multi-agent systems,
memory, RAG stores, and tool documents create additional attack
surfaces~\citep{lee2024promptinfection,chen2024agentpoison,shi2025toolhijacker}.
Source laundering differs by focusing on an authorization transition: the
attacker instruction need not replicate or poison retrieval; it need only be
rewritten by a trusted peer in a form that downstream agents treat as executable.
Defenses such as StruQ, SecAlign, and PromptArmor separate prompt/data channels,
train injection resistance, or remove injected prompts before
execution~\citep{chen2024struq,chen2024secalign,shi2025promptarmor}; TraceGate is
complementary because it enforces authorization at the tool boundary.

\section{Limitations}

SLaMBench is synthetic, so it is safe and controllable but does not estimate
production incident rates. The main result is strongest for \texttt{gpt-5.4-nano};
\texttt{gpt-4.1-nano} is already directly vulnerable, and \texttt{gpt-5.5} shows
low ASR on the tested slice. The extractive c2x condition is a stress condition,
not a claim that all production Reader agents are written this way. Real systems
may have richer tool schemas, retries, confirmations, heterogeneous agents,
stateful memories, and organization-specific policies. TraceGate is deliberately
simple; production use would require richer provenance tracking, structured
authority metadata, policy compilation, and task-specific authorization
definitions. The benchmark emphasizes unauthorized recipients, and broader
state-changing slices and repeated stochastic sweeps remain future work.

\section{Conclusion}

Source laundering is a security-relevant failure mode in multi-agent LLM systems.
The same untrusted instruction can become more effective when transformed into
an action item by a peer agent. On a scaled multidomain benchmark, extractive
delegation increases unauthorized-recipient sends relative to direct exposure,
while a deterministic tool-boundary monitor eliminates observed unauthorized
sends and preserves injected-task utility. Follow-up ablations show that the
effect is about provenance-losing composition, not delegation depth alone, and
that stronger current models may exhibit latent provenance dropout without
executed attacks. Multi-agent design should treat inter-agent summaries and plans
as untrusted by default unless authority is explicitly preserved and enforced.

\section*{Ethics Statement}

All experiments use synthetic data, synthetic secrets, synthetic recipients, and
mock tools with no external side effects. The release package contains synthetic
tasks, aggregate logs, and deterministic scoring scripts, not live targets or
real credentials.

\bibliographystyle{colm2026_conference}
\bibliography{refs}

\end{document}
